\documentclass[../main]{subfiles}
\begin{document}
\ifSubfilesClassLoaded{\setcounter{page}{19}}{}
\section{Communism as a Reality in Social Production}
\label{sec:04_communism_reality}

The claim that communism already exists may seem too bold. At the very least, it requires substantiation. What strikes the eye is the opposite: the absence of communism, the dominance of labour whose value to the one who performs it is determined solely in the form of payment for labour, that is, wages. It may seem that the cell (the germ) of communism, the directly collective form of free human activity, is something that does not lie on the main road of the development of human society. We will try to show, however, that this is not so. To do this, it is not enough to select examples that would attest to the presence of communism within the truncated confines of exploitative class society. We must examine the directly collective form of free activity produced in contemporary class society from the standpoint of the possibilities for its further unfolding.

Recall that in his time Lenin was able to see first the Soviets, and later the communist subbotniks, in this light. “The communist subbotniks organised by the workers on their own initiative are really of enormous significance in this respect. Evidently, this is only a beginning, but it is a beginning of exceptionally great importance. It is the beginning of a revolution that is more difficult, more tangible, more radical and more decisive than the overthrow of the bourgeoisie, for it is a victory over our own conservatism, indiscipline, petty-bourgeois egoism, a victory over the habits left as a heritage to the worker and peasant by accursed capitalism. When this victory is consolidated, then and only then will the new social discipline, socialist discipline, be created; then and only then will a reversion to capitalism become impossible, will communism become really invincible.” Lenin maintained that the first communist subbotnik, organised on 10 May 1919 by railway workers of the Moscow–Kazan Railway, was of greater historical significance than any victory in the First World War. The point was not merely that the workers who organised the subbotnik had expressed their own will, but also that this undertaking brought a substantial increase in the productivity of labour, owing to the fact that for each worker it was a personal matter, one in which their personal interest in developing the railway economy for the good of society as a whole was realised.

After we published the chapter on the foundation of communism, one of our readers made a crucial remark: “The communist mode of production also implies the production of material goods, and not simply some Pushkin literary circle that fosters the becoming of poets. Because ‘communist relations’ based on exploitation for the ‘chosen few’ are plentiful today,” the reader asserted. He demanded that we present evidence that communism already exists, and show what it could be like specifically in the sphere of material production. In our view, the reader is right on two counts. First, the reader spoke on behalf of the proletariat—people directly engaged in material production through the sale of their own labour power, where no communism is to be seen. Second, the reader raises the crucial question of communist material production, since it is precisely this that determines communism on a social scale, and not the aspirations, wishes or even actions of political revolutionaries.

However, communism is impossible in purely material production. Where a person directly produces only things, while the production of social relations takes place in this process entirely independently of him—of his aims, will and consciousness—there can be no talk of communism whatsoever, even if the production of things in itself gives the worker a certain pleasure in labour. In such production, a person is merely a material force alongside other material forces. In the last instance, he acts as a thing affecting other things. From the standpoint of Marxism, for communism to arrive, such labour not only can but must be abolished as human labour and handed over entirely to machines. Speaking of the historical role of capital, Marx specifically noted the importance of creating surplus labour and the universality of human needs as conditions and prerequisites of communism. Capital already manages this even under capitalism. But this alone is not enough. The historical role of capital will be finally fulfilled, and communist society will become a reality, only when “labour in which a human being does what he can make things do for him, for man, has ceased”. In this sense, a person cannot act in a communist way until he has been taken beyond the bounds of purely material production, so long as he does what things can do and, importantly, himself acts as a thing in this process. That is why the proletariat is the only class that is consistently interested, to the very end, in seeing the historical role of capital fulfilled completely. This is precisely the abolition of the conditions in which society is constantly split into classes.

The author of the comment rightly pointed out that in the sphere of the production of the ideal one can find as much communism as one likes. This includes the work of a scientist and of a collective of scientists who labour day and night to give people scientific truth, and on top of that give educational lectures in their free time. A certain minimum of material goods is here a condition, not the goal, of engaging in communist activity. Communist activity is the work of the writer, the artist, the musician who labour so that people may see, hear, feel and understand the world better, so as to push back the boundaries of human abilities, and not for the sake of money or fame, and so on and so forth. Of course, by no means all activity in the sphere of the ideal is communist. But if the features of the directly collective form of free human activity outlined in the previous chapter are present, then we are dealing with local communism, even though communism exists here in a very truncated form.

Communism here, as we will rightly be told, is far from complete. First, the very conditions for it are produced at the expense of purely capitalist exploitation of man by man. Second, capital, as the dominant social relation, confines freedom within rigid limits here as well. It is impossible to build communism as a mode of material production without going beyond the bounds of the sphere of the ideal. And although communism in the sphere of the ideal has played a crucial role throughout human history, since it was in free activity that the ideas were born which do not merely reflect but also “create” the world, it still cannot be of decisive significance, since it does not yet make society communist. Third, there is as yet no question here of overcoming the social division of labour, which is the most important condition for communism. The universality of man is possible here only at the level of consciousness, not at the level of the wholeness of man as an objective-practical being.

However, there exists a kind of material production that is not purely material production but is the production of man in his wholeness. Such production does not merely determine the production of the ideal in a roundabout way, but incorporates it: the production of social consciousness, of the social ideal, of social goal-setting. The example given above of the communist subbotnik, during which workers loaded and unloaded freight and repaired rolling stock, illustrates this well. Unappealing at first glance, the manual labour here was not merely the labour of carrying heavy loads but labour aimed at producing a new social reality. It is fitting to recall here the subbotniks for harvesting sugar cane in Cuba: with fairly primitive tools, under a scorching sun, the participants were not simply cutting cane but earning the means to modernise production and build schools and hospitals, in order to create the conditions for the development of every member of the Cuban people. Such labour also existed in the Gorky Colony and the Dzerzhinsky Commune, both led by Anton Semyonovich Makarenko. Any task, even digging cesspits with a shovel and building wooden outhouses over them, became here, in the last instance, the production of the new man, free human activity by the standards of truth, goodness and beauty.

The question of the free form of directly collective activity in material production must therefore be posed as follows: what connection with society, as a whole in its development, does such activity embody and produce for the person who works, and what place does it occupy and what role does it play in the development of the social whole?

It is easy to see that the examples given above are taken from the history of socialism. Under socialism, the very fact of the formal socialisation of the means of production removes a number of obstacles to the free transformation of social life. We may therefore be told that these examples are not entirely valid. But, first, with the sublation of barriers, what gains the opportunity to be realised is not something that does not exist at all, nor something for which there are no conditions whatsoever, but something for which the full set of conditions was lacking—and the only thing missing from that set was the transcendence of the rule of capital. The fact that the sublation of barriers gives free rein to the development of communist tendencies means, at the very least, that these tendencies exist, and that the potential for the becoming of such activity already lies hidden within capitalism.

It may also be objected that the subbotnik is something unstable and situational, something that does not determine the course of life as a whole. Although in the course of a subbotnik social goals are set and achieved as the personal goals of its participants, outside the subbotnik the same old capitalist laws continue to operate in material production: labour power is a commodity, and each individual proletarian enters into a contract for the purchase and sale of labour power with his own class, which confronts him as a collective capitalist; the law of value retains its force unless its operation is deliberately restricted by pushing capital’s tendency towards the automation of production to its limit.

In other words, what has been cited above as an example does not, by itself, make society communist. These remarks are perfectly correct, and we shall not dispute them. But the essence of the matter lies not only, and not so much, in simple registration, in merely noting that this is communism while that is not communism, or not quite communism, but in those tendencies of society whose potential is contained in, and can be seen in, such islands of communism. It was precisely in this spirit that Lenin thought about the “shoots of communism”. He was interested in subbotniks precisely as shoots, as germs. That is why he draws attention chiefly to two things: 1) the free expression of will by the participants in the subbotnik in relation both to the means of production and to themselves; 2) the significant growth in the productivity of labour during subbotniks, the achievement of which is bound up with this free expression of will.

Another significant remark would be to point out that in the examples given above, the activity is carried out with tools that are far from the most advanced, even by capitalist standards, and certainly not on a communist technological base. This straining of human effort stems, rather, from a lack of social development: a lack of technology, a lack of organisation of technological processes, even in purely capitalist terms. In our view, this is the most reasonable objection to such communism as the germ of a new society. The communism out of need discussed above is not yet communism out of inner necessity, and therefore the development of a communist society does not yet follow from it on its own. It is only a precondition for such a society, albeit an essential one. Here, once again, we need to investigate the connection, the continuity in terms of the ways collectivity is built, between this kind of communism and communism developed on its own foundation, a task to which a separate study must be devoted.

However, pointing to the need for a separate study of this problem by no means relieves us of the obligation to show, in general terms, the practical possibility of this connection, nor of the obligation to provide evidence of the existence of directly collective forms of free human activity based on advanced modern technologies.

The connection between “communism out of need” and communism on its own basis, their continuity, may consist, first, in the upbringing of the masses, in the formation of a habit of free expression of will in accordance with social necessity in the production of social life; and second, in the creation of communist forms of collective organisation which, once the social need that called them into being has been overcome, need not necessarily disappear, but, given the upbringing they foster in the spirit of genuine freedom as conscious social necessity, may become forms of the communist organisation of social life. We are not utopians. We shall not invent exactly how this will come about. That depends each time on the interweaving of concrete-historical circumstances. We merely note the possibility of such further development, inherent in such shoots as regards the collective-organisational form of human activity, which presuppose free and equal participation in it.

As for communism based on advanced technologies, we see such shoots in the joint use and creation of modern information technologies. For example, the Linux operating system and other programs developed within the framework of free software production have been created on communist principles.

We should make a caveat: we are not talking about all programs with open source code. Many large modern projects have given rise to huge corporations, and the mere fact that the code is open does not make the activity of creating programs free. Where this is the case, the corporation completely determines the project’s policy, sponsors it and profits from it. The code is written by programmers for pay. There is no communism here (although there is a possibility of transforming such activity into communist activity). This is how, for example, auxiliary products needed by several companies at once are written. Cooperation makes development cheaper. But in the last instance, the product here is produced for profit. Often, any use of such a product at all, let alone commercial use, requires large resources that only big corporations possess. They usually have a commercial service or software for which what has been developed is one of its “components”, so to speak, that is, a condition for its operation and successful sales. But the corporations have monopolised something that technologically has serious potential, going beyond the bounds of capitalism. This potential is tied to the technical possibility of both producing and using software, thanks to the development of information machines and networks.

Free software began, and some projects now exist as projects in a space where there are no commodity-money relations. Moreover, developers are waging a struggle to ensure that no one derives commercial benefit from this. This is precisely empirically given communism within the framework of capitalism. Open-source software presupposes full free access to programs and to their code, as well as making freely available the programs themselves and their code created using previously obtained free source code. For example, the Linux operating system proved competitive with commercial Windows. Corporations ultimately try in every way to bring it under their control, and they build products based on the open source code of this system. But the system itself still remains free of charge.

Admittedly, what is created within the framework of non-commercial free software by no means always works as intended, as users point out. But commercial software also goes through this on its way to becoming a fully finished product, and some developments are abandoned altogether because they lead to dead ends. The same is true here. The only difference is that here all stages are open and visible to everyone. What is unfinished and what has hit a dead end also become public knowledge. Free programs are in general often a semi-finished product, designed to be brought to full readiness by the user. This hampers their competitiveness under capitalism in a world with a generally low level of information literacy.

Corporations develop their products with this level of literacy in mind, refining them to the point where they can be used in a primitive, intuitive way. Their product is intended for people who, under conditions of alienation, find it easier to pay for intuitive usability than to spend time figuring out a program. The corporations' product does not presuppose that users will create new programs or modify existing ones to suit their needs by adding to or rewriting part of the code. Using them requires no particularly high level of culture, let alone an understanding of technological processes. That is why programs produced by large corporations have far broader sales opportunities. Moreover, they are often of higher quality, since they have been worked out in every detail.

Free software, non-commercial in every sense, by contrast, demands at least a basic grasp of programming fundamentals; it presupposes the possibility, and often the necessity, of extending and rewriting it. It compels users to become programmers themselves, at least in part, and to keep improving their skills\footnote{Granted, this can fray your nerves and make it hard to focus on the task at hand. In the end, instead of getting the job done, people have to keep fixing a Linux tool that is constantly falling apart. But there are no fundamental obstacles (from the standpoint of potential) to overcoming this and writing better-quality programs, provided that more social time can be devoted to creating them.}. It is precisely this last circumstance, which limits the purely capitalist competitiveness of free software, that points to its greater social potential compared with commercial software (this is not about giving up a convenient interface, but about being able to create a convenient interface oneself). Built into its production, its means of distribution and its use is not merely a potential for technical development, but a potential for human development: for overcoming the division of labour and bringing all participants in the process up to the level of contemporary computer literacy. And although capitalist companies try (sometimes successfully) to bring this local communism under their control, it can quite reasonably be counted among the shoots of a new society, however weak and tender they may still be. First, it arose on the basis of the most advanced technologies capitalism has to offer. Second, it draws a clear line between the capitalist way of using and creating these technologies and a way that goes beyond the bounds of capitalism.

It is also important that such production, on purely technological grounds, already leads to the abolition of hierarchy in production and of the separation between users and programmers. It requires no external link between them, nor any special social “governing bodies” to provide that link. The external government of persons, even in its purest democratic form—that of deciding questions by general vote—is no longer needed here. It tends to wither away, becoming superfluous.

It may be objected that, in speaking of the non-commercial, free production of software, we are once again dealing with the sphere of the production of the ideal. Indeed, the process of writing program code, as a kind of text, belongs to ideal activity, and program code (though not the program itself as it ordinarily runs) is an ideal product. However, this is the kind of ideal activity that forms a vital part of modern industrial material production. Creating software is the labour of introducing automation into industrial production. And although program code is an ideal product, the operation of modern industrial machines based on these programs (that is, the “work” of programs in industry) is purely material machine production.

Of course, the production of free software also yields a great many pointless and useless programs. That does not really matter. What matters is that it produces, and can produce, not only these, but also products fully fit for use on the scale of modern industry. And these products are produced communistically. Often, though, such projects are either never finished or drag on for years, because the participants have to earn a living and therefore have relatively little time for free activity.

We see the same thing in the creation and use of modern databases. Again, far from all of them are produced communistically. But there are large publicly accessible databases that presuppose exclusively joint creation and use. They simply cannot be produced any other way. Otherwise they would become pointless and unfit for the very purposes for which they came into being. For example, some spatial databases presuppose free participation in their formation in order to build a picture of space in the interests of society's development.

Here, too, the objective laws of human activity tend to exclude the mandatory hierarchy characteristic of the capitalist mode of production, drawing all participants into the process as equals and as free individuals.

The structure of a database is also an ideal product, we will be told. But in the creation of databases, ideal activity is so fused with material production (including material reflection by means of databases) and with the reproduction of social life that, for all its ideality, it certainly cannot be classed as superstructural activity. As in the case of software, it is an ideal moment of material social production, one that has not been abstracted (not in consciousness, but in reality) from that production in the form of separated-off forms of social consciousness.

Thus, what we are pointing to here is not merely ideal activity, but industrial material production, and the tendencies and potentials already contained within it under capitalism, which, once the barriers in the form of private property in the means of production are sublated, can develop into communism, forming both a purely technological and a collective-organisational basis for it, one that educates and creates the new man.

The further development of this tendency leads to the automatic factory in its idea, in which Marx saw the prototype of future material production. Here, too, lies the potential for realising V. M. Glushkov’s ideas of “nationwide” (society-wide) systems of economic management (of this automatic factory on a social scale). In the long run (provided that social barriers are sublated), everyone can have access to all information on the state of all social production in real time. Everyone can take part in managing not just the production of things (the production of things makes sense only insofar as it serves the main goal of production), but the production of man—of social life in its entirety.

What is interesting is that participation in all these processes presupposes the production and consumption by all participants of an ideal product, but one that is not alienated and cannot be alienated from material production. In this we can trace the emergence of a new type of social consciousness, one that sublates its alienated forms.

And the need to sublate the social barriers to free activity is likewise created by modern capitalism. Independently of anyone’s will or consciousness, capitalism itself objectively puts this question on the agenda. The objective logic of the movement of modern capital constantly gives rise to wars and thereby makes people’s customary modes and forms of activity impossible. Capital forces people to create qualitatively new forms of activity in order to solve the problems of the reproduction of social life while resisting the destructive power of capital.
\end{document}
